% Pursuing Stacks (cotes 134-2 to 134-8): archive cote 134-2, site page 74 (PDF page 75).
% appendix (letters to Breen). Draft compared: ps1.tex:1597-1614.
% AI-transcribed by Claude Opus 5.5 (claude-opus-5-5, adaptive thinking, effort high), Message Batches API, batch ps1 (msgbatch_01Rv5hCaNTtUt4JHx1yqoHPL), 2026-11-12; not yet checked by eye.
\scanpage{134-2/74}
\mspage{- \hand{5}9 -}

decomposition, i.e. the closed strata are again manifolds with boundary which
are glued together along common parts of the boundary; and it is also
necessary to consider the inclusions of these pieces one in\inserted{to} another ...],
and can be reconstituted from this diagram by gluing(\hand{50}).In other words, one
has a canonical devissage description of tame compact spaces $X$ , eventually
endowed with equi-singular stratifications with non-singular strata, in terms
of finite diagrams of a precise nature made out of manifolds with boundary.
When \deleted{in a situation of}\inserted{we are interested in} sheaves (or complexes of sheaves\hand{,} or n-stacks) which
are P-constructible on $X$ , where $P$ is such a fixed stratification, \deleted{known}\inserted{these}
\deleted{explicitly}\inserted{may be described\hand{,}} in terms of the envisaged diagram\hand{,}of which only the ``homotopy
type'' is to be retained. One \deleted{observes}\inserted{foresees} that the six operations on these sheaves
can be translated in an ad hoc manner to this homotopical context. Finally\hand{,}
if instead of having only one compact tame space $X$ , one has, let us say,
a tame morphism $f : X \to Y$ of such objects, then \deleted{by devissage of}\inserted{by choosing} equi-%
singular stratifications on $X$ and $Y$ adapted to $f$ (the strata of $X$
being in particular locally trivial fibrations on those of $Y$ ...), one
should find a ``morphism'' from the diagram of manifolds with boundary
expressing $X$ into that expressing $Y$ (with \deleted{some natural}\inserted{\unsure{mutal|natural}} morphisms which
essentially reduce to fibrations of \inserted{compact,} manifolds with boundaries on other such
objects) in such a way that the four operations $Rf_*$, $Rf_!$, $Lg^*$, $g^!$ between
$P_X$- and $P_Y$- constructible sheaves on $X$ and $Y$ (or on locally closed
sub spaces $X'$ , $Y'$ which are union of strata, such that $f$ induces
$g : X' \to Y'$) can be expressed in terms of standard operations between the
mere homotopy type\hand{s}. Finally, all these constructions, still partially
hypothetical (there is work on the foundations to be done!) should be able
to be paraphrased in the framework of excellent schemes, by making use of
the machinery of étale tubular neighbourhoods. In one or other case, the
``fine homotopy type'' of a tame space, respectively of an excellent scheme,
is defined by passage to the limit from ``P-homotopy type'' associated to finer
and finer equi-singular stratifications $P$ (with non-singular strata).

\margin{74}
